\textbf{Do we need to prove the existence of the `external' world?}

{\hfill\break
}\textbf{Abstract}

In this paper, I defend Wittgenstein's insight that Cartesian-style
skeptical arguments against the existence of the external world are not
false, but linguistic confusions that violate the constitutive rules of
our epistemic practices. To do this, I advance a resolutely grammatical
interpretation of his later notebooks, posthumously published as
\emph{On Certainty}, where so-called `hinge-propositions' function as
rules of grammar and logically exclude the possibility of `global'
doubt. I then defend this reading from two notable objections by Crispin
Wright, namely that (1) if `hinges' are not semantically truth-apt
propositions, they cannot function as premises in logical inferences,
and (2) the groundless nature of grammatical propositions renders the
foundation of our language-game arbitrary, leaving us unable to
distinguish a rational world-picture from mere dogma.

\textbf{Introduction}

Kant famously regarded the `scandal of philosophy' to be its inability
to present a sound proof for the existence of the world that could
withstand any doubt\footnote{Kant, \emph{Critique of Pure Reason},
  Bxxxix.}. While G.E. Moore argued that his proof accomplishes that
task, Wittgenstein maintained that the demand for a proof itself rested
on an incorrect understanding of language that, upon correction, would
render a proof unnecessary and logically exclude the possibility of
doubt. This paper proceeds in three main parts. In Section 1, I present
Wittgenstein's criticism of Moore's proof that it incorrectly uses the
term ``I know'', and argue that hinges function as grammatical
propositions. I also present a unique criticism of what I term P.M.S.
Hacker's `intermediate view' on the status of hinges. In Section 2, I
present Crispin Wright's theory of entitlement and warrant transmission
as a rival strategy before outlining two central objections he raises
against the grammatical reading. Finally, in Section 3, I develop a
novel defense against Wright's position and his specific objections. I
argue that treating hinges as rules of grammar does not leave our
epistemic practices at risk of arbitrariness or leaching off of unearned
warrant, which is a more robust dissolution of the problem of skepticism
than Wright's strategy. Furthermore, I establish that Wright's
entitlement theory functions as an epistemically inert pragmatism that
is excessively concessionary to the skeptic.

\textbf{1. Wittgenstein's Engagement with Moore}

\textbf{1.1 Moore's Proof and the Limits of Sense}

Wittgenstein's reflections respond primarily to G.E. Moore's two seminal
papers, `A Defence of Common Sense' (1925) and `Proof of an External
World' (1939).\footnote{G. E. Moore, \emph{Philosophical Papers} (1959).}
Moore sought to rebut idealism with an ostensibly basic proof by
declaring, while raising each hand in turn, ``Here is one hand'', and
``Here is another''. He insisted that this demonstration \emph{ipso
facto} proved an external world by meeting the formal requirements for
proof, namely, having epistemic certainty of the premises and logically
inferring the conclusion. Acknowledging the sceptic's objection
regarding unproven premises, Moore rejected the requirement for an
infinite regress of proof; instead, he maintained that certain
common-sense truths were known directly and had to be accepted to make
further reasoning possible.\footnote{Moore, \emph{Philosophical Papers},
  146--150.} Moore's strategy, later known as the ``Moorean
shift'',\footnote{See Rowe, \emph{"The Problem of Evil''} (1979),
  335--341.} was to reverse the sceptic's logic. While the sceptic used
\emph{modus tollens} to argue that a lack of knowledge regarding
skeptical hypotheses precludes knowledge of the external world, Moore
employed \emph{modus ponens} to argue that because he knew he had hands,
the skeptical hypotheses must be false.

Wittgenstein denied that Moore could claim knowledge of his premises,
arguing that the error lay in misusing the term ``I know.'' Central to
this critique is the logical requirement, tracing back to the
\emph{Tractatus,} that empirical propositions must be bipolar, i.e.
either true or false, to possess `sense' (TLP 4.022-4.023). Moore's
truisms lack this bipolar character\emph{,} prompting Wittgenstein to
label as ``obviously nonsensical'' any attempt to raise doubts where no
questions can be asked (TLP 6.51). This principle reappears later in
\emph{On Certainty,} where Wittgenstein maintains that knowledge-claims
are only intelligible where doubt is possible (OC 58). Because Moore's
assertions, unlike ordinary cases, cannot be justified by citing
evidence, they carry an illicit ``metaphysical emphasis'' on account of
their groundless nature (OC 482-484). They are more akin to a child's
acceptance of a ``tree'' through ostensive teaching (OC 480).
Consequently, these propositions function not as empirical descriptions
but as ``hinges'' (OC 341) that serve as the ``axis'' of all inquiry (OC
152-153). Just as a door requires fixed hinges to move, our empirical
inquiries depend on certain propositions to remain exempt from doubt so
that we can evaluate other claims. These propositions are termed hinges
to illustrate their critical role in our language-game. This marks an
anti-foundationalist shift that dissolves the infinite regress of
justification by recognizing the ``groundlessness of our believing'' (OC
166). If these hinges are not empirical, they must be understood as
constitutive rules of linguistic practice.

\textbf{1.2. Hinges as Grammatical Propositions}

In his transition from the \emph{Tractatus,} Wittgenstein reconceives
the status of foundational propositions, moving away from a strictly
picture theory of language. This meant that words and propositions only
possessed a sense if they functioned like a logical map that mirrored a
corresponding fact in the world. Language, then, was a descriptive
enterprise to `map' out all the facts of the world and form a picture.
However, in his later transitionary period after returning to Cambridge
from 1929 to 1936, Wittgenstein's understanding of language begins to
shift towards a rule-governed activity, as seen in his
\emph{Philosophical Grammar.} Wittgenstein recognizes that before any
`picturing' can occur, there must be a background of rules and
conventions that function as the conditions of picturing. This marks a
shift from a logic of representation to one of linguistic practice.
While maintaining that tautologies, contradictions, and logical rules
lack bipolarity (TLP 4.466), he departs from the standard
logico-semantic usage that reserves the term `proposition' exclusively
for truth-evaluable claims. This later shift marks a widening in his
understanding of grammar. Here, grammar is no longer merely a system of
syntax but encompasses ``all the conditions (the method) necessary for
comparing propositions with reality'' (PG I, 45).

Hinges, despite carrying the outward syntactic form of empirical
propositions, satisfy this deeper grammatical role. For example, the
sentence `nothing can be red and blue all over at the same time' isn't
discovered by examining thousands of objects to find a counter-example.
Instead, it defines how we use color-words. If someone were to claim
they found an object that was red and blue all over, we would question
if they understood the meaning of the words `red' and `blue' rather than
attributing a scientific discovery to their name. They function as
logical enabling conditions that make our epistemic practices possible
(OC 401-402). Consequently, I retain his technical use of `grammatical
proposition' {[}\emph{grammatischer Satz}{]} to capture their unique
logical function despite possessing the appearance of a descriptive fact
about the world. They act as `norms of representation' i.e. rules that
``describe the structure of language; describe its possibilities'' (BT,
149e). Because hinges are propositions that have been `hardened' from
what were once empirical observations (OC 96, 167), they are removed
from the realm of verification. They can no longer be tested like
hypotheses as they are the very criteria for what it means to `test'
anything at all. Acting as the framework of our language-game, they are
the `limiting cases' (TLP 4.466; cf. OC 54) that give our notions of
`truth' and `falsehood' their sense.

This constitutive role connects the establishment of a grammatical rule
directly to the adoption of a unit of measurement. By fixing the `unit',
grammar enables the use of empirical propositions to describe the world;
as Wittgenstein notes, while the rule itself establishes the metric, the
empirical proposition ``says how long an object is'' only by making use
of that established rule (BT, 189e). This illustrates how hinges, as
logical enabling conditions, do not describe facts but rather provide
the very possibility of factual description. Consequently, Wittgenstein
argues that grammar itself is autonomous and, being unaccountable to any
reality, is ultimately arbitrary (PG I, 133). Attempting to justify the
rules of grammar is a category mistake; it is nonsensical because they
define the terms of justification itself. As he notes: ``the conventions
of grammar can't be justified by a description of what is
represented\ldots{[}because{]} any description of that kind already
presupposes the rules of grammar'' (BT, 188e). Justification, therefore,
must come to an end. What lies at the bottom of the language-game is not
a self-evident truth that we perceive, but our way of acting. Hence,
grammatical propositions lack truth-aptness, as ``the ground is not
true, nor yet false'' (OC 204-205).

Because these foundational grounds are exempt from truth-evaluations,
their adoption is not compelled by objective facts but by our shared
practices. The `arbitrariness' of grammar is thus understood in this
strict sense of choosing a metric (BT, 186e). Just as the standard meter
bar in Paris (PI 50) is used to measure lengths but cannot itself be
measured against the scale it creates, hinges are used to `measure'
truth but cannot be subjected to those same evaluations. The logical
exclusion of doubt is absolute. To doubt a hinge is not to question a
fact, but to break the tool that makes its expression possible. Such
`global' doubt amounts to an ``annihilation of all yardsticks'' (OC
492). We recognize these hinges as the constitutive rules of our
language-game because of their peculiar role in logically excluding the
possibility of doubt (OC 194). This normative function of regulating the
boundaries of language, despite being `arbitrarily' chosen, is how
hinges ensure the stability of the entire framework of inquiry.

Given the unrefined nature of Wittgenstein's notebooks\footnote{See OC
  Preface, vii.}, the categorical status of hinges remains contested.
Two dominant interpretations emerge from this ambiguity\footnote{This
  account is not exhaustive. See Williams, \emph{Unnatural Doubts}
  (1991); Coliva, \emph{Moore and Wittgenstein}, (2010); Pritchard,
  \emph{Epistemic Angst} (2015).}. G.H. von Wright and Danièle
Moyal-Sharrock emphasize the 'animal' nature of these certainties,
situating them as non-propositional forms of praxis, which refers to the
unreflective and habitual ways of acting that underpin our
language-games.\footnote{See von Wright, \emph{Wittgenstein} (1982,
  178); Danièle Moyal-Sharrock, \emph{Understanding Wittgenstein}
  (2004).} Conversely, P.M.S. Hacker argues that they occupy a 'curious
intermediate position,' neither wholly empirical nor strictly
grammatical.\footnote{Hacker and Baker, \emph{Wittgenstein: Rules,
  Grammar and Necessity} (2009, 258).} He contends that they lack the
normative function required to determine concepts or inference
rules.\footnote{Hacker, \emph{Wittgenstein's Place in
  20\textsuperscript{th} Century Analytic Philosophy} (1996, 215).}
However, I argue that Hacker's refusal to grant hinges grammatical
status is difficult to reconcile with his own admission that empirical
propositions can be transformed into `norms of
representation'.\footnote{Hacker, \emph{Wittgenstein: Comparisons and
  Context} (2013, 161). For a detailed critique of Hacker's rejection of
  hinges as grammatical propositions, see Moyal-Sharrock, ``Beyond
  Hacker's Wittgenstein'' (2013).} It is important to note that while
Genia Schönbaumsfeld primarily defends the grammatical reading as a
rejoinder to Moyal-Sharrock's non-propositional account,\footnote{Schönbaumsfeld,
  ``Hinge Propositions and the ``Logical Exclusion of Doubt''\,'' (2017,
  104).} she does not explicitly extend this criticism to Hacker's
`intermediate' stance. Rather, Schönbaumsfeld's only mention of Hacker
is agreeing with him that only empirical propositions are bipolar while
grammatical remarks are not. She does not seem to note Hacker's position
is yet different, despite this brief congruence. I argue this only
further indicates that Hacker's rejection of hinges as grammatical
propositions is a dogmatic position, since he is already keen on
granting the possibility of other empirical propositions `hardening'
into rules, particularly in mathematics, that have the same normative
rule-like function as hinges. As such, I extend Schönbaumsfeld's
previously focused argument against Moyal-Sharrock into a broader view
that also extends to other competing interpretations of \emph{On
Certainty} and attempt to resolve previous contradictions into a clearly
defined position.

\textbf{1.3. Grammatical Analysis of the `External' World}

A consequence of the grammatical reading is the exposure of the
sceptic's error in bifurcating reality into `internal' and `external'
domains. Since linguistic meaning is tethered to a shared world, any
expression of doubt necessarily concedes the world's existence as a
condition of its own intelligibility. The sceptic's position is
therefore semantically nonsensical, as it attempts to negate the very
framework that makes their speech possible. As Wittgenstein notes, ``If
you are not certain of any fact, you cannot be certain of the meaning of
your words either'' (OC 114). On the reading I propose, hinges function
as logical enabling conditions that regulate the language-game rather
than describing a reality beyond it. While they enable the sensible
application of truth and falsehood, they are not themselves truth-apt in
an empirical sense. This reading, however, faces two significant
challenges from Crispin Wright's account of `epistemic entitlement': (1)
that the non-truth-aptness of hinges precludes them from serving as
premises in valid logical inferences, thereby blocking the
warrant-transmission required to justify our ordinary beliefs, and (2)
that the autonomy of grammar poses an epistemic challenge of
arbitrariness, threatening the rationality of our foundational
commitments.

{\hfill\break
}

\textbf{2. Hinges without Warrant?}

\textbf{2.1. Wright's Diagnosis: Transmission Failure}

Wright develops his critique by reconstructing Moore's argument with a
tripartite schematic of Type-I, Type-II, and Type-III propositions.
Here, (I) the subjective premise `I seem to have hands', provides
warrant for (II) the empirical claim that `a hand exists', which in turn
entails (III) the metaphysical conclusion that `an external world
exists'. Type-I propositions describe sensory experiences that justify
empirical Type-II propositions about the external world, such as Moore's
truisms. However, warrant can only transmit from Type-I to Type-II
propositions if we already possess antecedent warrant for Type-III
propositions. He defines these `cornerstones' as foundational
commitments such that, if we lack warrant for them, we cannot claim
justification for the set of empirical propositions, or as Wright puts
it, a \emph{`region of thought'}, that depends upon them for their
validity. \footnote{Wright, "Warrant for Nothing", 167; henceforth WN.
  See also Davies, "Epistemic Entitlement," 219n2, who observes a
  nuanced distinction between not having rational warrant and not being
  able to claim it.}.

Moore's premises, then, are ``information-dependent'' on Type-III
propositions because sensory experiences can only ground empirical
claims if we already grant causal interaction with the world. Since
belief in an external world lacks antecedent evidential warrant,
skeptical scenarios remain live possibilities that threaten to undercut
the evidence of our senses. As such, Wright concludes that Moore's proof
is epistemically circular.\footnote{Wright, "Wittgensteinian
  Certainties," 23, 26--28; henceforth WC. Wright, WN, 171--172.}

Wright tentatively draws support for his position from Wittgenstein's
own observations that the grounds for common-sense propositions are
never ``as certain as the very thing they were supposed to be grounds
for''. Since the grounds are no surer than the assertions themselves,
such propositions are instead affirmed ``without special testing'' (OC
136, 243, 307). Nevertheless, Wright fundamentally diverges from reading
hinges as non-truth apt grammatical propositions by reintroducing a need
for rational justification that he holds can be acquired
non-evidentially.

\textbf{2.2. Epistemic Entitlement}

Wright correctly recognizes that Moore's proof cannot avoid transmission
failure if evidential justification is the only way to acquire warrant
for cornerstones. To avoid this impasse, Wright proposes a ``unified
strategy'' that expands the definition and scope of acquiring warrant to
disjunctively encompass evidential justification, based on proof, and
entitlement, based on trust. In other words, Wright now argues that a
cornerstone can be rationally held even in the absence of evidence,
provided it meets the criteria for entitlement (WN 177, 208).

Wright introduces the concept of `entitlement' to \emph{accept}, rather
than \emph{believe}, any proposition P despite lacking evidential
warrant. He concedes that the sceptic is correct to claim that we lack
justified \emph{belief} in cornerstones but argues this does not prevent
our \emph{acceptance} of them (WN 174-176). Whereas belief is a doxastic
attitude requiring evidence and conviction of truth, acceptance is the
practical attitude of acting on the assumption that P. Unlike justified
belief, entitlement is a rational trust in accepting P. Wright divides
this non-evidential warrant into various categories. First, Wright
constructs a dominant strategy he terms `Strategic Entitlement' and
analogizes it to the decision-theoretic model in Hans Reichenbach's
defense of induction. Reichenbach agreed with Hume that \emph{epistemic}
justification for induction was impossible to prove through both \emph{a
priori} and \emph{a posteriori} means but argued that we are rationally
entitled to using it because it is the only decision that offers a
possibility of success, whereas skepticism guarantees failure (WN
178-180). Likewise, Wright's strategy reframes Reichenbach's pragmatic
argument to a form of epistemic consequentialism. It justifies the
acceptance of cornerstones based on the epistemic consequence of
securing the possibility of cognitive inquiry. Second, Wright identifies
the `Entitlement to Cognitive Project' as an entitlement to
presuppositions that stave off cognitive paralysis from an infinite
regress of doubt.

Wright insists that this appeal to entitlement does not constitute a
mere dogmatic dismissal of the sceptic's challenge\footnote{Cf. OC 71,
  155, 498, 558, 559. These passages may imply dogmatism if hinges are
  treated as empirical propositions.} but is grounded in the
\emph{epistemic} necessity of trust. He notes that while these
entitlements are effective against Cartesian and Humean forms of
skepticism, they do not readily extend to ontological cornerstones (WN
188, 194-197). Nevertheless, because our lives as rational agents depend
on trusting Type-III propositions, Wright argues that avoiding them
makes life impossible since rational agency is a ``not an optional
aspect of our lives''. As such, we are justified in presupposing their
truth without needing to earn a warrant (WC 53).

Wright immediately identifies a serious problem that threatens to
undermine his project. Since Wright affirms that cornerstones are held
by trust, they are beyond evidential justification and introduce a
degree of epistemic risk. This unearned warrant might `leach up' and
contaminate the warrants for otherwise ordinary empirical claims, such
as visually perceiving one's hand. By highlighting the ``leaching
problem'', the skeptic can now argue that ordinary justified beliefs are
taken up on mere faith. To resolve this, Wright proceeds by stratifying
knowledge claims and restricting warrant to second-order claims. We
cannot claim to know that cornerstones are true, but we can still
possess first-order knowledge of them (WN 209). On this view, possession
of knowledge is ``liberal'', requiring no antecedent warrant, while
claiming it is ``conservative''.\footnote{Wright, "On Epistemic
  Entitlement (II)," 217--218; henceforth WSE.} Wright claims that
leaching would thus be confined to second-order warrants, i.e. claims of
possession, leaving our first-order warrants untroubled.

\textbf{2.3. The Semantic and Epistemic Challenge}

Wright's notion of epistemic entitlement presents a semantic challenge
to the grammatical reading of hinges. Namely, Wright questions how
hinges can enable us to distinguish between true and false beliefs if
their own semantic content is not truth-conditional. Wright maintains
that because empirical inquiry has an external objective, the
``divination of what is true and the avoidance of what is false'' (WC
43), our cornerstones must be answerable to this truth-seeking goal. The
rules of arithmetic, like `13+7=20', must not force us to ``misassess
the truth-values'' of the specific things we are investigating. Wright
presents an example of a person counting 13 satsumas and 7 bananas but
mistakenly arriving at a count of 21 fruits, compelling the person to
conclude that they miscalculated. Here, the rule overrides the sensory
evidence because its certainty is ranked higher. However, if hinges were
merely groundless norms of representation, they would have no authority
to dismiss the physical possibility that there really are 21 fruits.
This is because, for Wright, unless a rule like 13+7=20 is truth-apt, it
cannot determine or describe what empirically exists in the world.
Wright concludes that if a rule is to reliably guide us to the truth, it
must be a ``substantial proposition, apt to be true or false'' (WC 34,
44-46). Denying this would leave cornerstones incapable of functioning
as premises in valid inferences, like Moore's proof, and `saves'
skepticism only by expelling them from the realm of logic entirely.

Even if cornerstones, such as the Earth having existed for many years
past (OC 138, 401), serve as rules, Wright argues this is not mutually
exclusive with their being truth-evaluable. Interpreting cornerstones as
\emph{hors de combat} i.e., incapable of intelligent mismatch with the
world, is hasty and does not banish ``the spectre of profound and
sweeping error'' (WC 44).

In addition to the semantic challenge, Wright identifies a deeper
epistemic threat of arbitrariness that looms over the grammatical
reading. If grammar is autonomous and, by Wittgenstein's own admission
(PG 133), arbitrary, it's difficult to see how we can secure a rational
standing for epistemic practices. Although Wright acknowledges
cornerstones are themselves unfounded {[}\emph{unbegründet}{]} by
calling them presuppositions (2004b, 48-49) he still maintains the need
for a rational basis in affirming and acting upon them. Wright terms the
``Problem of Demarcation'' as the danger of being unable to distinguish
cornerstones from ``prejudices, mere assumptions, and idées fixes'' (WSE
245). If our cornerstones lack rational warrant, we risk succumbing to a
pernicious relativism where we cannot distinguish a legitimate
constitutive rule and a mere subjective whim.

To avoid running into an infinite regress of demands for justifications,
Wright's tentative response is to reject what he calls the ``Evidential
Ideal'', the notion that a rationally epistemic thinker should
proportion all their beliefs to evidence and accept nothing without it.
Wright argues that the ideal itself is incoherent because evidential
reasons to believe can only occur within a framework held by trusting
acceptances. This trust, far from being a flaw, is an essential aspect
of rational inquiry without which we lose the ability to determine what
counts as evidence (WSE 242-245). For example, determining why
perceptual experiences relate to matters in the external world
necessitates the acceptance of cornerstone that ``There is an external
material world, broadly manifest in normal sensory experience'' (WSE
217). The grammatical reading, despite claiming logical immunity from
skepticism, does not concord with the life of a rational agent precisely
because it is free from the necessity of trusting cornerstones. As such,
for Wright this reading fails because it cannot distinguish itself from
the arbitrary acceptance of dogma.

\textbf{3. Response to Wright}

\textbf{3.1. The Logical Barrier}

Wright is correct to identify that if hinges are not truth-apt, then
they cannot function as premises in any logical inferences, including
Moore's proofs. Moreover, he is right to note that their non-descriptive
nature, as norms of representation lacking sense, precludes them from
being items of knowledge.

However, I reject that this has any epistemically catastrophic
consequences. First, Wright's rejection of the grammatical reading rests
on a subtle category mistake that conflates the rules of the
language-game with the moves made within it. Because Wright treats the
foundations of our rational system and the empirical claims made within
it as belonging to the same class of propositions, namely empirical
ones, he incorrectly demands that they be subject to the same epistemic
evaluations. To follow Wittgenstein's analogy of language as a game of
chess, Wright's approach is akin to claiming that the rule `The bishop
moves diagonally' is of the same kind as the report 'The white bishop
moved from its starting square at c1 to f4'. The first sentence only
expresses a rule that defines how the bishop functions in every game
whereas the latter describes a specific instance of play. Wright
overlooks the notion that the scope of rationality, so to speak, is
`local'; it can only take place against a backdrop of propositions, our
world-picture as the inherited background, that we don't stand in any
epistemic relation to. Hinges are, as Schönbaumfeld notes,
`\emph{logical enabling conditions'} that give our notions of truth and
falsity their sense and without which we would lose the very meaning of
our words. Their scaffolding status serves to clarify the correct use of
language. Moore's very attempt to use hinges, along with his belief that
proof was necessary, was mistaken precisely because he treated the
constitutive rules of grammar as if they were empirical claims subject
to evaluation.\footnote{Genia Schönbaumsfeld, "`Logical' and `Epistemic'
  Uses," 239, 244--245.}

In contrast to Wright's assertion that hinges must be empirical
propositions in need of warrant, a resolute grammatical reading
maintains that this criterion is unnecessary because it is logically
impossible to doubt hinges. Wittgenstein repeatedly analogizes the rules
of grammar to the rules of chess (BT, 118e) to illustrate that
questioning them is no more sensible than questioning why the bishop
moves diagonally.

Wright's confusion lies in the fact that he views cornerstones to be
descriptive facts about reality. But rather than describing empirical
facts, the meaning of grammatical propositions lies in their use in
language where their ``explanation of meaning'' (PG I, 23) serves a
prescriptive function; to establish the rules of our language-game and
`police', so to speak, the correct method of using concepts. This is the
key insight behind Wittgenstein's assertion that the proposition,
``there are physical concepts'', is senseless and cannot be formulated
(OC 36). Because it is impossible to make sense of the proposition's
denial, i.e. that there are only sense-experiences, it follows
\emph{ipso} \emph{facto} that the assertion itself is mistaken. Such
propositions find their only use as instructions. It does not follow,
therefore, from the lack of truth-aptness of hinges, and their
inadmissibility as premises that we cannot make any logical inferences
with empirical propositions. Rather, as I have already attempted to
explain, Moore's mistake lay in using hinges in his proof, for the exact
reason that Wright purports is allegedly a defect in the grammatical
reading.

Having shown that a well-defined distinction exists between hinges as
the rules of our language-game and the empirical propositions as `moves'
within it, I argue it becomes clear that Wright's concern about leaching
is dissolved. At the bottom of our language-game lie hinges that do not
provide warrant but logically enable the conditioning of sense. Since
`risk' is an epistemic concept, it cannot apply to the grammatical
framework itself. Thus, leaching does not seep up, and any attempt at
grounding this world-picture would only invite further grounding,
\emph{ad} \emph{infinitum}. In a recent debate with Annalisa Coliva and
Anil Gupta\footnote{See Coliva, Gupta, and Wright, "A Debate on
  Skepticism and Perceptual Belief," 361, 379.}, despite amending his
prior definition of entitlement by adding that it is the attitude
towards a cornerstone proposition P, Wright's view is still at risk of
such regress, for the rational trust he places in P remains exposed to
the demand of entitlement towards \emph{that} attitude as well. Wright's
demand for hinges to be truth-apt is an attempt to repair a defect that
does not exist if they are taken to constitute a logical barrier for our
language-game. By conceding that cornerstones require justification,
Wright leaves himself vulnerable to an infinite regress of justification
in addition to leaching. I argue this conundrum is downstream of the
category mistake I have identified in Wright's work so far. By treating
empirical claims and cornerstones as the same kind of proposition,
Wright renders the task of preventing leaching impossible.

\textbf{3.2. Contra Entitlement}

Wright explicitly rejects the label of pragmatist by insisting that
cornerstones are essential for the possibility of our cognitive
projects,\footnote{See Pritchard, "Wittgenstein's On
  Certainty,"189--224; Jenkins, "Entitlement and Rationality," 25--45.
  Cf. Wright, WSE, 213--248 for his defense against such
  "instrumentalist" interpretations.} and that this constitutes a
normative condition. Wright's decision-theoretic strategy contends that
because these entitlements aim to enhance ``expected epistemic utility''
within a cognitive project, they are also epistemic reasons (WSE 239,
241-242). However, I argue this appeal to utility and the goal-dependent
nature of the `dominant strategy' to avoid intellectual stasis
functionally represents a pragmatist solution despite Wright's
insistence to the contrary. Wright shifts the justification of
cornerstones from the evidence of their truth to the utility of their
acceptance as the means of attaining warrant, drawing on an earlier
contrast he makes between acceptance and beliefs (WN 175, 184-188). As
Wright acknowledges, ``there is no coherent conception of epistemic
rationality which views unevidenced trusting as per se a lapse'' (WSE
243). In other words, Wright's unified strategy remains an instrumental
form of rationality. But, as Wittgenstein is keen to remind us, the
sceptic's challenge is not one of practical doubt but a more radical
`global' form that lies behind it (OC 19). The sceptic would succeed in
pinpointing the presuppositions of our epistemic practices, and
consequently the very meaning of notions like ``justification'' and
``warrant'' that Wright is keen to employ.

Furthermore, Wright's stratification of warrant fails to dispel concerns
of leaching because, as Aidan McGlynn argues, Wright's liberalism
towards first-order warrants, such as having a justified belief in
Moore's premise (``Here is a hand'') in combination with the recognition
that it entails belief in the cornerstone proposition (``There is an
external world'') leads him to adopting a form of neo-Mooreanism that
delivers the antecedent warrant that Wright's conservatism at the
second-order demands, resulting in his entitlement theory being rendered
``an idle wheel''.\footnote{McGlynn, "Epistemic Entitlement and the
  Leaching Problem," 95--96, 98.}

Finally, I argue that the notion of an entitlement is epistemically
inert; it satisfies the structural requirement for cornerstone
propositions without providing the truth-conducive grounding necessary
to stop the leaching of unearned warrant. This results in a system of
conditional rationality where the cognitive project of our empirical
beliefs is internally consistent but anchored by a set of ``blind''
wagers that we are incapable of confirming. This underlying pragmatism
is notable in Wright's rejection of the ``Evidential Ideal'' (WSE 243)
on the grounds that it is an integral part of our rational reflective
inquiry that we take some items of knowledge on unearned trust, because
the alternative is ``intellectual stasis''. This, as I have attempted to
show, does not sufficiently address the `global' nature of the sceptic's
doubt.

\textbf{3.3. Hinges as Constitutive of Rationality}

This brings us to the core of Wright's `Problem of Demarcation': the
concern that treating hinges as groundless rules renders our
world-picture merely optional. I argue Wright's notion of `Strategic
Entitlement' treats cornerstones instrumentally as seen in the contrast
Wittgenstein makes between the rules of cooking and grammar; cooking is
defined by its end, whereas speaking a language isn't (PG I, 133). There
is a categorical difference between instrumental activities and
grammatical ones:

\begin{quote}
"I believe the reason one isn't tempted to call cooking a game is this:
... 'cooking' doesn't refer to an activity that follows those rules --
rather, to an activity that has a particular result... On the other hand
``playing chess'' doesn't mean an activity that has a particular result,
but rather an activity that corresponds to such and such rules." (BT,
186e).
\end{quote}

If someone were to change or not follow the rules of cooking (e.g.,
boiling an egg for five minutes instead of three), he might fail to
achieve the intended result; but if we change the rules of grammar, we
don't speak ``falsely'', rather we are playing a different game. The
notion of not following a rule of grammar is nonsense. Wright treats
hinges instrumentally in his decision-theoretic strategy, like a cooking
recipe designed to achieve a desired outcome: truth or survival.
Secondly, Wittgenstein's assertion that grammar is arbitrary does not
imply hinges are mere prejudices, assumptions, or \emph{idées} fixes, as
Wright worries by identifying a ``Problem of Demarcation'' (WSE 245) but
only that one shouldn't confuse a rule with the use of words or expect
it to be answerable to reality in the way an empirical proposition is.
This is because grammar is autonomous; it does not `mirror' the world
but rather provides the framework that makes any description of the
world possible. The conventions of grammar cannot be justified by a
description of what is being represented, as ``any description of that
kind already presupposes the rules of grammar'' (BT, 188e).

As I established earlier, because hinges are logical enabling conditions
like the constitutive rules of a game,\footnote{Schönbaumsfeld,
  "`Logical' and `Epistemic' Uses," 241.} they cannot be random
assumptions that we could conceivably make differently, owing to their
contingent nature. The arbitrariness of grammatical rules stems from our
observing their fixity in language before removing these propositions
from the realm of contingency.

\begin{quote}
Consider: ``In language the only correlate to natural necessity is an
arbitrary rule. It is the only thing one can remove from this necessity
and put into a proposition.'' Refers to propositions such as
\textasciitilde{} \textasciitilde p = p. (BT, 185e)
\end{quote}

Just as the rule of double negation is not an empirical fact we
discover, but a constitutive rule informing us of how the word ``not''
is used, likewise the arbitrariness of grammar consists in the fact that
it is not responsible to any reality; it is not a description of facts,
but the very framework that makes descriptions of facts possible.
Furthermore, our world-picture of hinges are not axioms that provide the
starting point for language, but the very medium in which rationality
exists.

\begin{quote}
``And this system is not a more or less arbitrary and doubtful point of
departure for all our arguments: no, it belongs to the essence of what
we call an argument. The system is not so much the point of departure,
as the element in which arguments have their life.'' (OC 105)
\end{quote}

Wright's mistake lies in treating cornerstones as these points of
departure since they open us to question if alternate points of
departure could be chosen or not, so to speak. By insisting on the
empirical and contingent nature of cornerstones, Wright's entitlements
are vulnerable to the same problem of demarcation he identified earlier.

\textbf{4. Conclusion}

The debate between Wright and the grammatical reading I have proposed
turns on the nature of the ground of our epistemic practices. Wright
accepts the sceptic's demand that cornerstones require warrant and
offers `entitlement' as part of his unified strategy. Yet, as I have
argued in this paper, Wright's strategy is epistemically inert because
it doesn't adequately resolve the problem of leaching, leaving us with a
system that is only conditionally rational. By conceding to the skeptic
that our cornerstones require warrant, Wright leaves our entire
epistemic practices vulnerable to the leaching problem. Because
entitlement is mere trust, it cannot prevent the further contamination
of our empirical claims. In contrast, by acknowledging hinges as
grammatical propositions, I argue such doubts are logically excluded and
that the arbitrariness of grammar is constitutive, functioning like the
rules of a game that define the field of play rather than the moves made
within it. I provide a more robust dissolution of the skeptical
challenge by identifying Wright's category mistake through my analogies
with chess games. I show that hinges are the constitutive grammatical
rules that define the moves made within our language game. This is a
significant shift towards an anti-foundationalist attempt at securing
the rational integrity of our epistemic practices without having to
justify, potentially infinitely, each cornerstone proposition.

\textbf{List of Abbreviations}

\begin{longtable}[]{@{}ll@{}}
\toprule
\textbf{Abbreviation} & \textbf{Full Title \& Edition}\tabularnewline
\midrule
\endhead
\textbf{BT} & \emph{The Big Typescript: TS 213}, ed. and trans. C.G.
Luckhardt and M.A.E. Aue (Oxford: Blackwell, 2005).\tabularnewline
\textbf{OC} & \emph{On Certainty}, ed. G.E.M. Anscombe and G.H. von
Wright, trans. D. Paul and G.E.M. Anscombe (Oxford: Blackwell,
1969).\tabularnewline
\textbf{PG} & \emph{Philosophical Grammar}, ed. R. Rhees, trans. A.
Kenny (Oxford: Blackwell, 1974).\tabularnewline
\textbf{PI} & \emph{Philosophical Investigations}, trans. G.E.M.
Anscombe, P.M.S. Hacker, and J. Schulte, 4th ed. (Oxford:
Wiley-Blackwell, 2009).\tabularnewline
\textbf{TLP} & \emph{Tractatus Logico-Philosophicus}, trans. D.F. Pears
and B.F. McGuinness (London: Routledge \& Kegan Paul,
1961).\tabularnewline
\textbf{WC} & \emph{Wittgenstein and Skepticism,} ed. Denis McManus
(London: Routledge, 2004).\tabularnewline
\textbf{WN} & \emph{``Warrant for Nothing (and Foundations for Free)?''}
The Aristotelian Society Supp. 78, no.1 (2004)\tabularnewline
\textbf{WSE} & \emph{"On Epistemic Entitlement (II): Welfare State
Epistemology."} In Scepticism and Perceptual Justification (Oxford: OUP
2014)\tabularnewline
\bottomrule
\end{longtable}

\hypertarget{bibliography}{%
\section{Bibliography}\label{bibliography}}

Coliva, Annalisa. \emph{Moore and Wittgenstein: Scepticism, Certainty
and Common Sense.} Basingstoke, UK: Palgrave Macmillan, 2010.Coliva,
Annalisa, Anil Gupta, and Crispin Wright. "A Debate on Skepticism and
Perceptual Belief." In \emph{Empirical Reason and Sensory Experience},
edited by Ori Beck and Miloš Vuletić, 355--390. Cham, Switzerland:
Springer, 2024 https://doi.org/10.1007/978-3-031-52231-4\_45.Davies,
Martin. "Epistemic Entitlement, Warrant Transmission And Easy
Knowledge." \emph{Aristotelian Society Supplementary} 78, no. 1 (2004):
213-245. https://doi.org/10.1111/j.0309-7013.2004.00122.x.Hacker, P. M.
S. \emph{Wittgenstein's Place in Twentieth-Century Analytic Philosophy.}
Oxford: Blackwell, 1996.---. \emph{Wittgenstein: Comparisons and
Context.} Oxford: Oxford University Press, 2013.Hacker, P. M. S. , and
G.P. Baker. \emph{Wittgenstein: Rules, Grammar and Necessity: Essays and
Exegesis of 185--242.} Oxford: Blackwell, 2009.Jenkins, C. S.
"Entitlement and rationality." \emph{Synthese} 157, no. 1 (2007): 25-45.
https://doi.org/10.1007/s11229-006-0012-2.Kant, Immanuel. \emph{Critique
of Pure Reason.} Translated and edited by Paul Guyer and Allen W. Wood.
Cambridge: Cambridge University Press, 1998.McGlynn, Aidan. "Epistemic
Entitlement and the Leaching Problem." \emph{Episteme} 14, no. 1 (2017):
89-102. https://doi.org/10.1017/epi.2015.63.Moore, G.E.
\emph{Philosophical Papers.} London: Allen and Unwin,
1959.Moyal-Sharrock, Danièle. \emph{Understanding Wittgenstein.} New
York: Palgrave-Macmillan, 2004.---. "Beyond Hacker's Wittgenstein:
Discussion of HACKER, Peter (2012) ``Wittgenstein on Grammar, Theses and
Dogmatism'' Philosophical Investigations 35:1, January 2012, 1--17."
\emph{Philosophical Investigations} 36, no. 4 (2013): 355-380.
https://doi.org/10.1111/phin.12021.Pritchard, Duncan. "Wittgenstein's On
Certainty and Contemporary Anti-skepticism." In \emph{Readings of
Wittgenstein's On Certainty}, edited by Danièle Moyal-Sharrock and
William H. Brenner, 189-224. London: Palgrave Macmillan, 2005.---.
\emph{Epistemic Angst: Radical Skepticism and the Groundlessness of Our
Believing.} Princeton, NJ: Princeton University Press, 2015.Rowe,
William. "The Problem of Evil and Some Varieties of Atheism."
\emph{American Philosophical Quarterly} 16, no. 4 (1979): 335-341.
https://www.jstor.org/stable/20009775.Schönbaumsfeld, Genia. "`Hinge
Propositions' and the `Logical' Exclusion of Doubt." In \emph{Hinge
Epistemology}, edited by Annalisa Coliva and Danièle Moyal-Sharrock,
94--109. Leiden, Netherlands: Brill, 2017.
https://doi.org/10.1163/9789004332386\_007.---. "`Logical' and
`Epistemic' Uses of `to Know' or `Hinges' as Logical Enabling
Conditions." In \emph{Skeptical Invariantism Reconsidered}, by Genia
Schönbaumsfeld, edited by Christos Kyriacou and Kevin Wallbridge,
235-251. London: Routledge, 2021.Williams, Michael. \emph{Unnatural
Doubts: Epistemological Realism and the Basis of Scepticism.} Oxford:
Blackwell, 1991.Wittgenstein, Ludwig. \emph{Tractatus
Logico-Philosophicus.} Translated by D. F. Pears and B. F. McGuinness.
London: Routledge \& Kegan Paul, 1961.---. \emph{On Certainty.} Edited
by G. E. M. Anscombe and G. H. von Wright. Translated by Denis Paul and
G. E. M. Anscombe. Oxford: Blackwell, 1969.---. \emph{Philosophical
Grammar.} Edited by Rush Rhees. Translated by Anthony Kenny. Oxford:
Blackwell, 1974.---. \emph{The Big Typescript: TS 213.} Edited by C.
Grant Luckhardt and Maximilian A. E. Aue. Translated by C. Grant
Luckhardt and Maximilian A. E. Aue. Oxford: Blackwell, 2005.---.
\emph{Philosophical Investigations.} 4th. Translated by G. E. M.
Anscombe, P. M. S. Hacker and J. Schulte. Oxford: Wiley-Blackwell,
2009.Wright, Crispin. "Warrant for Nothing (and Foundations for
Free)?,." \emph{The Aristotelian Society} 78, no. 1 (2004): 167-212.
https://doi.org/10.1111/j.0309-7013.2004.00121.x.---. "Wittgensteinian
Certainties." In \emph{Wittgenstein and Skepticism}, by Crispin Wright,
edited by Denis McManus, 22--55. London: Routledge, 2004.---. "On
Epistemic Entitlement (II): Welfare State Epistemology." Chap. 10 in
\emph{Scepticism and Perceptual Justification}, edited by Elia Zardini
and Dylan Dodd. Oxford: Oxford University Press, 2014.Wright, Georg
Henrik von. \emph{Wittgenstein}. Minneapolis, MN: University of
Minnesota Press, 1982.
